\documentclass[12pt]{article}
\usepackage{latexblog}

% ============================================================
% Blog Metadata
% ============================================================
\blogtitle{Jer语言(1)：语法设计}
\blogdate{2021-01-05}
\blogtags{JVM, Bytecode, 闲话编程, Antlr, 小实现}
\bloglang{zh}
\blogsummary{}

\begin{document}
\makeblogtitle

最近准备实现一个基于JVM的新语言''Jer''，一开始想先实现一个Hello world，然后逐步再朝上面添加新的功能；后来觉得还是需要先把这个语言的语法层面大致设计好再动手才行。本身是出于好玩的一个目的，但是也的确希望这个语言有一些特点，而不是单纯换一个语法而已。在这个期间思考了很多，但一直没有想到自己满意的方法，姑且先按照现在的想法设计一版出来吧。

\section{设计目标}\label{ux8bbeux8ba1ux76eeux6807}

我对这个语言有这些期望：

\begin{itemize}
\tightlist
\item
  基于JVM平台，即最终通过代码编译生成class字节码
\item
  跟Java能够兼容（可以互相调用）
\item
  尽量简单，应该基本的数据类型、流程控制等，支持OO，但对于一些高级特性例如泛型、lambda等就不考虑了
\item
  依然是强类型的语言
\end{itemize}

\subsection{hello world}\label{hello-world}

类似Java这样一个java文件只能对应一个类也许不是也个好的办法，jer源文件中可以申明任意数量的类。对于一个hello world来说，(大概）应该长这样子：

\begin{Shaded}
\begin{Highlighting}[]
\CommentTok{// 导入其他类或者其中的静态方法}
\KeywordTok{use}\NormalTok{ jer}\OperatorTok{/}\NormalTok{lang}\OperatorTok{/}\NormalTok{System}

\CommentTok{// 没有定义在类中的方法对应到java中的静态方法}
\NormalTok{main(args}\OperatorTok{:}\NormalTok{ [}\DataTypeTok{String}\NormalTok{) }\OperatorTok{=} \OperatorTok{\{}
\NormalTok{    msg}\OperatorTok{:} \DataTypeTok{String} \OperatorTok{=} \StringTok{"hello world!"}
\NormalTok{    println(msg)}
\OperatorTok{\}}
\end{Highlighting}
\end{Shaded}

其中，不需要使用分号作为行的分隔符，直接换行就行了。一个Jer源文件即对应到一个Java的类，类的名称即是文件名。

\section{Jer 语法}\label{jer-ux8bedux6cd5}

\begin{Shaded}
\begin{Highlighting}[]
\NormalTok{compilationUint}
\NormalTok{    : importedType* declaration* EOF}
\NormalTok{    ;}
\NormalTok{declaration}
\NormalTok{    : constantDeclaration}
\NormalTok{    | methodDeclaration}
\NormalTok{    | abstractDeclaration}
\NormalTok{    | typeDeclaration}
\NormalTok{    ;}
\end{Highlighting}
\end{Shaded}

每一个文件中，可以包含任意：

\begin{itemize}
\tightlist
\item
  导入申明
\item
  常量变量申明
\item
  方法申明（即为静态函数）
\item
  type或者类的申明
\end{itemize}

\subsection{导入申明}\label{ux5bfcux5165ux7533ux660e}

通过导入申明来引入其他包或者文件中定义的类。

\begin{Shaded}
\begin{Highlighting}[]
\NormalTok{importedType}
\NormalTok{    : USE fullPath}
\NormalTok{    ;}
\NormalTok{fullPath}
\NormalTok{    : (IDENTIFIER \textquotesingle{}/\textquotesingle{})* TYPE\_NAME}
\NormalTok{    ;}
\end{Highlighting}
\end{Shaded}

\begin{Shaded}
\begin{Highlighting}[]
\KeywordTok{use}\NormalTok{ java}\OperatorTok{/}\NormalTok{lang}\OperatorTok{/}\DataTypeTok{String}
\KeywordTok{use}\NormalTok{ java}\OperatorTok{/}\NormalTok{util}\OperatorTok{/}\NormalTok{DateTime}
\end{Highlighting}
\end{Shaded}

这里需要考虑一个场景就是，如果是其他Jer文件中定义了常量或者静态方法，在其他文件中如何使用？

\begin{Shaded}
\begin{Highlighting}[]
\CommentTok{// com/riguz/jer/Util.jer}
\NormalTok{sum(a}\OperatorTok{:}\NormalTok{ Integer}\OperatorTok{,}\NormalTok{ b}\OperatorTok{:}\NormalTok{ Integer) }\OperatorTok{{-}\textgreater{}}\NormalTok{ Integer }\OperatorTok{=} \OperatorTok{\{}
    \CommentTok{// ...}
\OperatorTok{\}}

\CommentTok{// com/riguz/jer/Foo.jer}
\KeywordTok{use}\NormalTok{ com}\OperatorTok{/}\NormalTok{riguz}\OperatorTok{/}\NormalTok{jer}\OperatorTok{/}\NormalTok{Util}

\NormalTok{main(args}\OperatorTok{:}\NormalTok{ [}\DataTypeTok{String}\NormalTok{) }\OperatorTok{=} \OperatorTok{\{}
\NormalTok{    sum(}\DecValTok{1}\OperatorTok{,} \DecValTok{20}\NormalTok{)}
\OperatorTok{\}}
\end{Highlighting}
\end{Shaded}

\subsection{常量定义}\label{ux5e38ux91cfux5b9aux4e49}

常量定义跟普通的局部变量唯一的区别就是多了一个const的关键字。具体的语法在后面介绍。

\begin{Shaded}
\begin{Highlighting}[]
\NormalTok{constantDeclaration}
\NormalTok{    : CONST variableDeclaration}
\NormalTok{    ;}
\end{Highlighting}
\end{Shaded}

\subsection{方法}\label{ux65b9ux6cd5}

\begin{Shaded}
\begin{Highlighting}[]
\NormalTok{methodDeclaration}
\NormalTok{    : methodSignature methodImplementation?}
\NormalTok{    ;}
\NormalTok{methodSignature}
\NormalTok{    : IDENTIFIER \textquotesingle{}(\textquotesingle{} formalParameters? \textquotesingle{})\textquotesingle{} functionReturnType?}
\NormalTok{    ;}
\NormalTok{formalParameters}
\NormalTok{    : formalParameter (\textquotesingle{},\textquotesingle{} formalParameter)*}
\NormalTok{    ;}
\NormalTok{functionReturnType}
\NormalTok{    : TO type}
\NormalTok{    ;}
\NormalTok{methodImplementation}
\NormalTok{    : \textquotesingle{}=\textquotesingle{} block}
\NormalTok{    ;}
\NormalTok{formalParameter}
\NormalTok{    : IDENTIFIER \textquotesingle{}:\textquotesingle{} type}
\NormalTok{    ;}
\end{Highlighting}
\end{Shaded}

方法分为两种，一种是有返回值的方法，另一种是没有返回值的方法（void），在定义的时候稍微有些区别：

\begin{Shaded}
\begin{Highlighting}[]
\CommentTok{// 没有返回值的方法依靠方法的副作用}
\NormalTok{main(args}\OperatorTok{:}\NormalTok{ [}\DataTypeTok{String}\NormalTok{) }\OperatorTok{=} \OperatorTok{\{}
\NormalTok{    msg}\OperatorTok{:} \DataTypeTok{String} \OperatorTok{=} \StringTok{"hello world!"}
\NormalTok{    println(msg)}
\OperatorTok{\}}

\CommentTok{// 返回值通过箭头表示}
\NormalTok{sum(a}\OperatorTok{:}\NormalTok{ Integer}\OperatorTok{,}\NormalTok{ b}\OperatorTok{:}\NormalTok{ Integer) }\OperatorTok{{-}\textgreater{}}\NormalTok{ Integer }\OperatorTok{=} \OperatorTok{\{}
    \ControlFlowTok{return}\NormalTok{ a }\OperatorTok{+}\NormalTok{ b}
\OperatorTok{\}}
\end{Highlighting}
\end{Shaded}

\subsection{抽象类和自定义类型}\label{ux62bdux8c61ux7c7bux548cux81eaux5b9aux4e49ux7c7bux578b}

\begin{Shaded}
\begin{Highlighting}[]
\NormalTok{abstractDeclaration}
\NormalTok{    : ABSTRACT TYPE\_NAME \textquotesingle{}\{\textquotesingle{} propertyDeclaration* methodSignature*\textquotesingle{}\}\textquotesingle{}}
\NormalTok{    ;}
\NormalTok{typeDeclaration}
\NormalTok{    : TYPE TYPE\_NAME typeAbstractions? \textquotesingle{}\{\textquotesingle{} propertyDeclaration* constructorDeclaration* methodDeclaration*\textquotesingle{}\}\textquotesingle{}}
\NormalTok{    ;}
\NormalTok{typeAbstractions}
\NormalTok{    : IS TYPE\_NAME (\textquotesingle{},\textquotesingle{} TYPE\_NAME)*}
\NormalTok{    ;}
\NormalTok{propertyDeclaration}
\NormalTok{    : IDENTIFIER \textquotesingle{}:\textquotesingle{} type}
\NormalTok{    ;}
\NormalTok{constructorDeclaration}
\NormalTok{    : \textquotesingle{}(\textquotesingle{} constructorFormalArguments? \textquotesingle{})\textquotesingle{} methodImplementation}
\NormalTok{    ;}
\NormalTok{constructorFormalArguments}
\NormalTok{    : constructorFormalArgument (\textquotesingle{},\textquotesingle{} constructorFormalArgument)*}
\NormalTok{    ;}
\NormalTok{constructorFormalArgument}
\NormalTok{    : IDENTIFIER (\textquotesingle{}:\textquotesingle{} TYPE\_NAME)?}
\NormalTok{    ;}
\end{Highlighting}
\end{Shaded}

\subsection{数据类型}\label{ux6570ux636eux7c7bux578b}

\begin{Shaded}
\begin{Highlighting}[]
\NormalTok{type}
\NormalTok{    : TYPE\_NAME}
\NormalTok{    | arrayType}
\NormalTok{    ;}
\NormalTok{arrayType}
\NormalTok{    : \textquotesingle{}[\textquotesingle{} type}
\NormalTok{    ;}
\end{Highlighting}
\end{Shaded}

\subsubsection{基本数据类型}\label{ux57faux672cux6570ux636eux7c7bux578b}

数据类型与Java基本一致，对应到JVM的各个数据类型：

\begin{itemize}
\tightlist
\item
  Bool : JVM boolean
\item
  Byte : JVM byte
\item
  Short: JVM short
\item
  Integer: JVM int
\item
  Long: JVM long
\item
  Float: JVM float
\item
  Double: JVM double
\item
  Char: JVM char
\item
  String: java/lang/String
\end{itemize}

取消java中的primitive 类型，即所有一切都是引用类型。

\subsubsection{数组类型}\label{ux6570ux7ec4ux7c7bux578b}

数组类型用\texttt{{[}\textless{}Type\textgreater{}}表示，例如\texttt{{[}Integer}即表示一个整数数组。

\subsection{表达式}\label{ux8868ux8fbeux5f0f}

\begin{Shaded}
\begin{Highlighting}[]
\NormalTok{expression}
\NormalTok{    : primary}
\NormalTok{    | expression bop=\textquotesingle{}.\textquotesingle{}}
\NormalTok{        ( methodCall}
\NormalTok{        | IDENTIFIER}
\NormalTok{        )}
\NormalTok{    | methodCall}
\NormalTok{    | objectCreation}
\NormalTok{    ;}
\NormalTok{primary}
\NormalTok{    : \textquotesingle{}(\textquotesingle{} expression \textquotesingle{})\textquotesingle{}}
\NormalTok{    | literal}
\NormalTok{    | IDENTIFIER}
\NormalTok{    ;}
\NormalTok{literal}
\NormalTok{    : DECIMAL\_LITERAL}
\NormalTok{    | FLOAT\_LITERAL}
\NormalTok{    | CHAR\_LITERAL}
\NormalTok{    | STRING\_LITERAL}
\NormalTok{    | BOOL\_LITERAL}
\NormalTok{    | NULL\_LITERAL}
\NormalTok{    ;}
\NormalTok{methodCall}
\NormalTok{    : instance=IDENTIFIER? \textquotesingle{}(\textquotesingle{}methodName=IDENTIFIER methodArguments? \textquotesingle{})\textquotesingle{}}
\NormalTok{    ;}
\NormalTok{methodArguments}
\NormalTok{    : expression (\textquotesingle{},\textquotesingle{} expression)*}
\NormalTok{    ;}
\NormalTok{objectCreation}
\NormalTok{    : NEW \textquotesingle{}(\textquotesingle{} methodArguments? \textquotesingle{})\textquotesingle{}}
\NormalTok{    ;}
\end{Highlighting}
\end{Shaded}

\subsection{statement}\label{statement}

\begin{Shaded}
\begin{Highlighting}[]
\NormalTok{block}
\NormalTok{    : \textquotesingle{}\{\textquotesingle{} statement* \textquotesingle{}\}\textquotesingle{}}
\NormalTok{    ;}
\NormalTok{statement}
\NormalTok{    : variableDeclaration}
\NormalTok{    | embeddedStatement}
\NormalTok{    ;}

\NormalTok{embeddedStatement}
\NormalTok{    : block}
\NormalTok{    | assignment}
\NormalTok{    | expressionStatement}
\NormalTok{    | selectionStatement}
\NormalTok{    | loopStatement}
\NormalTok{    | returnStatement}
\NormalTok{    ;}
\NormalTok{assignment}
\NormalTok{    : IDENTIFIER \textquotesingle{}=\textquotesingle{} expression}
\NormalTok{    ;}
\NormalTok{selectionStatement}
\NormalTok{    : IF \textquotesingle{}(\textquotesingle{} expression \textquotesingle{})\textquotesingle{} statement (ELSE statement)?}
\NormalTok{    ;}
\NormalTok{loopStatement}
\NormalTok{    : WHILE \textquotesingle{}(\textquotesingle{} expression \textquotesingle{})\textquotesingle{} statement}
\NormalTok{    ;}
\NormalTok{returnStatement}
\NormalTok{    : RETURN expression}
\NormalTok{    ;}
\NormalTok{expressionStatement}
\NormalTok{    : methodCall}
\NormalTok{    ;}
\NormalTok{variableDeclaration}
\NormalTok{    : IDENTIFIER \textquotesingle{}:\textquotesingle{} type (\textquotesingle{}=\textquotesingle{} variableInitializer)?}
\NormalTok{    ;}
\NormalTok{variableInitializer}
\NormalTok{    : arrayInitializer}
\NormalTok{    | expression}
\NormalTok{    ;}
\NormalTok{arrayInitializer}
\NormalTok{    : \textquotesingle{}\{\textquotesingle{} variableInitializer (\textquotesingle{},\textquotesingle{} variableInitializer)* \textquotesingle{}\}\textquotesingle{}}
\end{Highlighting}
\end{Shaded}

\section{代码示例}\label{ux4ee3ux7801ux793aux4f8b}

\begin{Shaded}
\begin{Highlighting}[]
\KeywordTok{use}\NormalTok{ java}\OperatorTok{/}\NormalTok{lang}\OperatorTok{/}\DataTypeTok{String}
\KeywordTok{use}\NormalTok{ java}\OperatorTok{/}\NormalTok{util}\OperatorTok{/}\NormalTok{DateTime}

\KeywordTok{const}\NormalTok{ pi}\OperatorTok{:}\NormalTok{ Float }\OperatorTok{=} \DecValTok{3}\OperatorTok{.}\NormalTok{1415926f}
\KeywordTok{const}\NormalTok{ msg}\OperatorTok{:} \DataTypeTok{String} \OperatorTok{=} \StringTok{"hello world"}
\KeywordTok{const}\NormalTok{ kb}\OperatorTok{:}\NormalTok{ Integer }\OperatorTok{=} \DecValTok{1024}
\KeywordTok{const}\NormalTok{ success}\OperatorTok{:}\NormalTok{ Boolean }\OperatorTok{=} \ConstantTok{true}
\KeywordTok{const}\NormalTok{ id}\OperatorTok{:}\NormalTok{ Long }\OperatorTok{=} \DecValTok{12345678}

\NormalTok{main(args}\OperatorTok{:}\NormalTok{ [}\DataTypeTok{String}\NormalTok{) }\OperatorTok{=} \OperatorTok{\{}
\NormalTok{    (println }\StringTok{"Hello world"}\NormalTok{)}
\OperatorTok{\}}

\NormalTok{circleArea(radius}\OperatorTok{:}\NormalTok{ Float) }\OperatorTok{{-}\textgreater{}}\NormalTok{ Float }\OperatorTok{=} \OperatorTok{\{}
    \ControlFlowTok{return}\NormalTok{ pi(multiply }\DecValTok{2}\OperatorTok{,}\NormalTok{ radius)}
\OperatorTok{\}}

\KeywordTok{abstract}\NormalTok{ Movable }\OperatorTok{\{}
\NormalTok{    x}\OperatorTok{:}\NormalTok{ Integer}
\NormalTok{    y}\OperatorTok{:}\NormalTok{ Integer}
    \KeywordTok{move}\NormalTok{(x1}\OperatorTok{:}\NormalTok{ Integer}\OperatorTok{,}\NormalTok{ y1}\OperatorTok{:}\NormalTok{ Integer)}
\OperatorTok{\}}

\KeywordTok{abstract}\NormalTok{ Animal }\OperatorTok{\{}
\NormalTok{    name}\OperatorTok{:} \DataTypeTok{String}
\NormalTok{    sayHelloTo(person}\OperatorTok{:}\NormalTok{ Person)}
\NormalTok{    address() }\OperatorTok{{-}\textgreater{}} \DataTypeTok{String}
\OperatorTok{\}}

\KeywordTok{type}\NormalTok{ Dog is Animal, Movable \{}
\NormalTok{    (name) }\OperatorTok{=} \OperatorTok{\{}
\NormalTok{        x }\OperatorTok{=} \DecValTok{0}
\NormalTok{        y }\OperatorTok{=} \DecValTok{0}
    \OperatorTok{\}}

\NormalTok{    sayHelloTo(person}\OperatorTok{:}\NormalTok{ Person) }\OperatorTok{=} \OperatorTok{\{}
\NormalTok{        (println }\StringTok{"Hello"}\NormalTok{)}
    \OperatorTok{\}}
\OperatorTok{\}}
\end{Highlighting}
\end{Shaded}


\end{document}
